===== BEGIN FILE: article.tex =====

\documentclass[sn-mathphys,Numbered]{sn-jnl}


\usepackage{graphicx}
\usepackage{multirow}
\usepackage{amsmath,amssymb,amsfonts}
\usepackage{amsthm}
\usepackage{mathrsfs}
\usepackage[title]{appendix}
\usepackage{xcolor}
\usepackage{textcomp}
\usepackage{manyfoot}
\usepackage{booktabs}
\usepackage{algorithm}
\usepackage{algorithmicx}
\usepackage{algpseudocode}
\usepackage{listings}


\usepackage{subcaption}
\usepackage[section]{placeins}
\usepackage{multirow}

\usepackage{comment}



\newcommand\aap{A\&A}
\let\astap=\aap
\newcommand\aapr{A\&ARv}
\newcommand\aaps{A\&AS}
\newcommand\actaa{Acta Astron.}
\newcommand\afz{Afz}
\newcommand\aj{AJ}
\newcommand\ao{Appl. Opt.}
\let\applopt=\ao
\newcommand\aplett{Astrophys.~Lett.}
\newcommand\apj{ApJ}
\newcommand\apjl{ApJ}
\let\apjlett=\apjl
\newcommand\apjs{ApJS}
\let\apjsupp=\apjs


\newcommand\apss{Ap\&SS}
\newcommand\araa{ARA\&A}
\newcommand\arep{Astron. Rep.}
\newcommand\aspc{ASP Conf. Ser.}
\newcommand\azh{Azh}
\newcommand\baas{BAAS}
\newcommand\bac{Bull. Astron. Inst. Czechoslovakia}
\newcommand\bain{Bull. Astron. Inst. Netherlands}
\newcommand\caa{Chinese Astron. Astrophys.}
\newcommand\cjaa{Chinese J.~Astron. Astrophys.}
\newcommand\fcp{Fundamentals Cosmic Phys.}
\newcommand\gca{Geochimica Cosmochimica Acta}
\newcommand\grl{Geophys. Res. Lett.}
\newcommand\iaucirc{IAU~Circ.}
\newcommand\icarus{Icarus}
\newcommand\japa{J.~Astrophys. Astron.}
\newcommand\jcap{J.~Cosmology Astropart. Phys.}
\newcommand\jcp{J.~Chem.~Phys.}
\newcommand\jgr{J.~Geophys.~Res.}
\newcommand\jqsrt{J.~Quant. Spectrosc. Radiative Transfer}
\newcommand\jrasc{J.~R.~Astron. Soc. Canada}
\newcommand\memras{Mem.~RAS}
\newcommand\memsai{Mem. Soc. Astron. Italiana}
\newcommand\mnassa{MNASSA}
\newcommand\mnras{MNRAS}
\newcommand\na{New~Astron.}
\newcommand\nar{New~Astron.~Rev.}
\newcommand\nat{Nature}
\newcommand\nphysa{Nuclear Phys.~A}
\newcommand\pra{Phys. Rev.~A}
\newcommand\prb{Phys. Rev.~B}
\newcommand\prc{Phys. Rev.~C}
\newcommand\prd{Phys. Rev.~D}
\newcommand\pre{Phys. Rev.~E}
\newcommand\prl{Phys. Rev.~Lett.}
\newcommand\pasa{Publ. Astron. Soc. Australia}
\newcommand\pasp{PASP}
\newcommand\pasj{PASJ}
\newcommand\physrep{Phys.~Rep.}
\newcommand\physscr{Phys.~Scr.}
\newcommand\planss{Planet. Space~Sci.}
\newcommand\procspie{Proc.~SPIE}
\newcommand\rmxaa{Rev. Mex. Astron. Astrofis.}
\newcommand\qjras{QJRAS}
\newcommand\sci{Science}
\newcommand\skytel{Sky \& Telesc.}
\newcommand\solphys{Sol.~Phys.}
\newcommand\sovast{Soviet~Ast.}
\newcommand\ssr{Space Sci. Rev.}
\newcommand\zap{Z.~Astrophys.}




\theoremstyle{thmstyleone}
\newtheorem{theorem}{Theorem}


\newtheorem{proposition}[theorem]{Proposition}


\theoremstyle{thmstyletwo}
\newtheorem{example}{Example}
\newtheorem{remark}{Remark}

\theoremstyle{thmstylethree}
\newtheorem{definition}{Definition}

\raggedbottom


\begin{document}



\title[Article Title]{Planet-induced Periodic Modulation of Stellar Activity in GJ 436: Insights into a Warm Neptune’s Magnetic Field}



\author*[1]{\fnm{D.} \sur{Revilla}}\email{drevilla@iaa.csic.es}

\author[1]{\fnm{P.J} \sur{Amado}}

\author[1,2,19]{\fnm{R.} \sur{Luque}}

\author[1]{\fnm{P.} \sur{Schöfer}}

\author[3]{\fnm{A.} \sur{Binnenfeld}}



\author[4]{\fnm{J.A.} \sur{Caballero}}



\author[5]{\fnm{Artie P.} \sur{Hatzes}}




\author[6]{\fnm{G. W.} \sur{Henry}}




\author[7]{\fnm{S.} \sur{Jeffers}}


\author[8,9]{\fnm{S.} \sur{Kaur}}


\author[10]{\fnm{A.F.} \sur{Lanza}}




\author[11,12]{\fnm{E.} \sur{Pallé}}


\author[1]{\fnm{L.} \sur{Peña-Moñino}}


\author[1,13]{\fnm{M.} \sur{Pérez-Torres}}


\author[14]{\fnm{A.} \sur{Quirrenbach}}


\author[15]{\fnm{A.} \sur{Reiners}}


\author[8,9]{\fnm{I.} \sur{Ribas}}


\author[8,9,16]{\fnm{D.} \sur{Viganò}}



\author[3,17]{\fnm{S.} \sur{Zucker}}












\affil*[1]{\orgname{Instituto de Astrofísica de Andalucía (IAA-CSIC)}, \orgaddress{\street{Glorieta de la Astronomía s/n}, \city{Granada}, \postcode{18008}, \country{Spain}}}

\affil[2]{\orgdiv{Department of Astronomy and Astrophysics}, \orgname{University of Chicago}, \orgaddress{\street{5640 South Ellis Avenue}, \city{Chicago}, \postcode{60637}, \state{IL}, \country{USA}}}

\affil[3]{\orgname{Porter School of the Environment and Earth Sciences}, \orgdiv{Raymond and Beverly Sackler Faculty of Exact Sciences, Tel Aviv University}, \orgaddress{Tel Aviv, \postcode{6997801}, \country{Israel}}}

\affil[4]{\orgdiv{Centro de Astrobiología (CAB), CSIC-INTA, ESAC Campus}, \orgaddress{\street{Camino bajo del castillo s/n},\city{Villanueva de la Cañada, Madrid}, \postcode{28692}, \country{Spain}}}

\affil[5]{\orgdiv{Thüringer Landessternwarte Tautenburg}, \orgaddress{\street{Sternwarte 5},\city{Tautenburg}, \postcode{07778}, \country{Germany}}}

\affil[6]{\orgdiv{Tennessee State University (retired)}, \orgaddress{\city{Nashville}, \postcode{ 37209}, \country{USA}}}

\affil[7]{\orgname{Max Planck Institute for Solar System Research}, \orgaddress{\street{Justus-von-Liebig-Weg 3}, \city{Göttingen} \postcode{D-37077}, \country{Germany}}}

\affil[8]{\orgname{Institut d’Estudis Espacials de Catalunya (IEEC)}, \orgaddress{\street{Carrer d'Esteve Terradas, 1},\city{Castelldefels, Barcelona}, \postcode{08860}, \country{Spain}}}

\affil[9]{\orgname{Institut de Ciències de I'Espai (CSIC-IEEC)}, \orgdiv{Campus UAB}, \orgaddress{\street{Carrer de Can Magrans s/n},\city{Cerdanyola del Vallès, Barcelona}, \postcode{08193}, \country{Spain}}}

\affil[10]{\orgdiv{INAF-Osservatorio Astrofisico di Catania}, \orgaddress{\street{Via S. Sofia, 78}, \city{Catania}, \postcode{95123}, \country{Italy}}}

\affil[11]{\orgname{Instituto de Astrofísica de Canarias (IAC)}, \orgaddress{ \city{La Laguna, Tenerife} \postcode{38205}, \country{Spain}}}

\affil[12]{\orgname{Universidad de La Laguna (ULL)}, \orgdiv{Departamento de Astrofísica}, \orgaddress{\street{Diogenes street}, \city{La Laguna, Tenerife} \postcode{38206}, \country{Spain}}}

\affil[13]{\orgname{School of Sciences}, \orgdiv{European University Cyprus}, \orgaddress{\street{Diogenes street}, \city{Nicosia, Engomi} \postcode{1516}, \country{Cyprus}}}

\affil[14]{\orgname{Landessternwarte}, \orgdiv{ Zentrum für Astronomie der Universität Heidelberg}, \orgaddress{\street{Königstuhl 12},\city{Heidelberg}, \postcode{69117}, \country{Germany}}}

\affil[15]{\orgname{Institut für Astrophysik und Geophysik}, \orgdiv{Georg-August-Universität}, \orgaddress{\street{Friedrich-Hund-Platz 1},\city{Göttingen}, \postcode{37077}, \country{Germany}}}

\affil[16]{\orgname{Institute of Applied Computing \& Community Code (IAC3)}, \orgdiv{University of the Balearic Islands}, \orgaddress{\city{Palma}, \postcode{07122}, \country{Spain}}}

\affil[17]{\orgname{School of Physics and Astronomy}, \orgdiv{Raymond and Beverly Sackler Faculty of Exact Sciences, Tel Aviv University}, \orgaddress{Tel Aviv, \postcode{6997801}, \country{Israel}}}

\affil[18]{\orgname{Institute of Applied Computing \& Community Code (IAC3)}, \orgdiv{University of the Balearic Islands}, \orgaddress{\city{Palma}, \postcode{07122}, \country{Spain}}}

\affil[19]{NHFP Sagan Fellow}


\abstract{\textbf{Interactions between stellar and planetary magnetic fields are expected to produce observable radio and optical signals modulated by their orbital periods, but direct detections remain elusive. We analyze 17 years of spectroscopic data of the GJ 436 system. This M2.5 V star hosts a transiting Neptune-sized planet in a close-in, inclined orbit. The data shows repeated enhancements of the stellar chromospheric activity at approximately the same phase of its 8-year activity cycle modulated by a combination of the planet’s orbital period and the stellar rotation. We interpret this modulation as star-planet interaction. We propose a new geometrical model to interpret these signals, then, estimate the power of the interaction and, from models, estimate the magnetic field of GJ 436 b to be between 6 and 110G. This finding opens new pathways to detect star-planet interactions and to investigate planetary magnetic fields and their implications on atmospheric retention and detectability.}}

\keywords{Exoplanet systems, Planetary magnetic fields, Star-planet interactions, Stellar magnetic fields}


\maketitle

\section*{}\label{sec1}


GJ 436 ($d=9.78\, \textrm{pc}$, $R=0.42 \, R_{\odot}$, $M=0.44 \, M_{\odot}$ and $P_{\rm rot}=45\pm5 \, \textrm{d}$, \cite{Rosenthal_2021}, \cite{Bourrier_2018}) is an M2.5V type star hosting a transiting warm Neptune in a close-in orbit of 2.64\,d (\cite{Butler_2004}). Due to its favorable observability characteristics: close-in orbit, planet size, and bright host, the system is an ideal target for producing and investigating magnetic star-planet interactions (SPI). Prior studies have focused on the planet’s extended, escaping atmosphere (\cite{Bourrier_2016}) and its interaction with the stellar wind (\cite{Bourrier_2016},\cite{Cauley_2017},\cite{Lavie_2017}), suggesting potential for detectable SPI. However, no confirmation of direct observation of magnetic SPI has been reported for this or any other system to date.  In this work, we analyze 371 high-resolution spectra of GJ~436 taken with the instruments HARPS (\cite{Mayor_2003}) (170 observations between 2006 and 2010, plus 23 observations in 2020) and CARMENES (127 observations in 2016 and 51 observations in 2024) (\cite{Quirrenbach_2014}). HIRES (\cite{Vogt_1994}) also observed this system but the S/N of the data is too low to provide results in our analysis (see Methods for more details). \\

We searched for periodic signals in several stellar activity indicators covered by the HARPS and CARMENES spectrographs using a suite of tools designed to detect time-varying and transient signals (see Methods). The best results were obtained with the Generalized Lomb-Scargle periodogram (GLS) (\cite{Zechmeister_2009}) in combination with rolling periodograms (\cite{Herbort_2018},\cite{Schofer_2019}) (see details in Methods).
We found two significant periods at 2.81\,d (and its corresponding 1-day alias at 1.54\,d) in HARPS 2008 and CARMENES 2016 datasets (hereinafter H-08 and C-16), and at 2.46\,d (and its corresponding 1-day alias at 1.68\,d) in CARMENES 2024 dataset (C-24) (Fig. \ref{Fig_1}). These periodicities are symmetric around the planet’s orbital period in frequency space and can be described by the relation $P_{\text{SPI}\pm} = (P_{\text{orb}}^{-1} \pm P_{\text{rot}}^{-1})^{-1}$, with $P_{\text{SPI-}}$ corresponding to the synodic period, i.e., the time it takes for a body to complete one rotation relative to the object it orbits. The signals appeared predominantly in the CaII~{\textsc{IRT-a}} and CaII~{\textsc{H\&K}} lines, which are strongly correlated (\cite{Busa_2007}, \cite{Martin_2017}) and are well-established tracers of chromospheric activity through their respective indices, pEW'(CaII~{\textsc{IRT-a}}) (\cite{Schofer_2019}) and $I_{CaII}$ (\cite{Kumar_2022}). We analyzed all other available activity indicators both from the photosphere and the chromosphere (see Methods) and found no further indication of the presence of these or any other signals, indicating that the emission is located at a certain altitude in the stellar chromosphere (\cite{Hintz_2023}). \\


To explain the detected signals, we propose a geometrical model with two independent active regions: one related to the stellar activity and another one generated by the SPI. The intrinsic stellar activity would produce the first, while the second would be the consequence of GJ~436~b interacting with GJ 436 through the magnetic field. Therefore one would move at the stellar rotation period, whereas the other at the orbital period. As a consequence of the extra energy coming from the interaction at the footpoint of those magnetic field lines, an enhancement of the chromospheric activity is produced in that region (hot spot). The model considers the rotation period of the star, the orbital period of the planet, and also the obliquity of the planetary orbit (see details in Methods). We find that, when the obliquity is zero, $I=0$, the only period that appears is $P_{\text{SPI}-}=P_{\text{syn}}$. However, for an oblique orbit, as in the case of GJ~436~b with $I=103\pm13^{\text{o}}$ (\cite{Bourrier_2018}), the other period, $P_{\text{SPI}+}$, emerges, and a beat-like signal is produced as seen in the synthetic data generated by our model in Fig. \ref{Fig_2}. Furthermore, we incorporated the phase difference between the position of the two active regions in the model. Depending on that phase difference, we can either detect the modulation at $P_{\text{SPI}-}$, $P_{\text{SPI}+}$ or sometimes both periods in the periodogram (Fig. \ref{Fig_2}, see details in Methods). From Eq. \eqref{dist2} and the synthetic data, we can see that there is no combination of phases that would produce a signal at $P_{\rm orb}$, even when $I=0$. \\


Our geometrical model cannot explain why the chromospheric stellar activity was only enhanced in 2008, 2016, and 2024. M~dwarfs are known to have activity cycles with periods of the order of years, appearing in both the photometry and in the chromospheric indices (\cite{Vida_2013},\cite{Robertson_2013},\cite{Suarez-Mascareño_2016},\cite{Küker_2019}). Previous studies have estimated the activity cycle of GJ 436 to be approximately $7.75 \pm 0.10$ years (\cite{Lothringer_2018, Loyd_2023}). We reanalyze the archival photometry together with new data from the T12 Automatic Photoelectric Telescope (APT) and the C14 Automated Imaging Telescope (AIT) after 2017 (see details in Methods) and show them in the context of the chromospheric calcium indices from HARPS and CARMENES in Fig.~\ref{Fig_3}. The new AIT observations show a less smooth and cyclic behavior but are still consistent with a variability of approximately 8\,years. The period of the activity cycle coincides with the appearance of the SPI signals in the spectroscopic data (2008, 2016, and 2024), suggesting that they occur approximately at a transitional phase between the activity cycle's maxima and minima. In addition, the total power produced by the SPI is directly related to the star's magnetic field strength (see Sec.~\ref{Sec:SPI-models}). Therefore, temporal changes in the strength of the star's magnetic field should directly influence the detectability of the SPI signal. While the maximum of the stellar activity cycle may seem optimal for detecting SPI due to the expected increase in power, larger active regions are expected during this time (\cite{Hathaway_2015}, \cite{Menezes_2023}), potentially diluting the signal from the SPI hot spot. On the other hand, the minimum of the activity cycle would correspond to a lower total SPI power output, making it very difficult to detect it. Our observations suggest that in this intermediate state, the magnetic field may be sufficiently strong to produce detectable SPI while the overall background activity level is lower, reducing the number of active regions that could otherwise dilute the signal.\\

The repeated appearance of signals in chromospheric activity indicators, at two distinct periods ($P_{\text{SPI}\pm}$) linked to the star and planet, and separated by approximately one stellar activity cycle (8 years), aligns well with our geometrical model and interpreting them as SPI. For detecting magnetic SPI, the planet must orbit within the Alfvén surface of the star --- the boundary where the star’s magnetic field dominates over forces such as stellar wind pressure. Spectropolarimetric observations obtained with NARVAL in 2016 \cite{Kumar_2022} allowed for the reconstruction of the star's magnetic topology and Alfvén surface \cite{Bellotti_2023, Vidotto_2023}, confirming that GJ~436~b resides within this region for most of its orbit during the C-16 observations. Given that H-08 and C-24 are separated by nearly one full activity cycle, we infer that the planet remained within the Alfvén surface during those epochs as well, where we detect additional SPI signals. Additionally, the signals were observed in the calcium lines, which have been the main tracers in the optical for SPI investigations \cite{Shkolnik_2003, Cauley_2018}.
Finally, the observed SPI periods are directly related to both the planet and the star. $P_{\text{SPI-}}$ is the synodic period $P_{\text{syn}}$ of the system, commonly predicted by SPI models to show signals of these interactions (\cite{Fischer_2019}). The second period, $P_{\text{SPI+}}$, arises particularly for this system due to the planet's highly inclined orbit.\\


Under this SPI interpretation, we can estimate the magnetic field of the planet, which is one of the hardest parameters to constrain observationally in exoplanet research. The first step is to estimate the total power released in the SPI. We used the $\chi$ factor calibrations (\cite{Walkowicz_2004},\cite{Labarga_2022}) to convert the pseudo equivalent width of the CaII~{\textsc{IRT-a}} line from C-16 and C-24 observations to total emitted power in the interaction, obtaining in both epochs powers around $10^{19}\,\text{W}$ (see details in Methods).
However, these detections represent only a fraction of the total power dissipated, as a larger fraction of power is expected to be radiated in wavelengths outside the calcium lines. To estimate the total power, we compared the SPI emission to the decay phase of a small flare in the M~dwarf AD Leo observed in multiple wavelengths (\cite{Hawley_2003}, \cite{Loyd_2023}) (for more details see Sec. \ref{Sec:Frac_Power}). Approximately 2\% of the total power is emitted in the CaII~{\textsc{IRT-a}} line during the decay phase, while half of it is in CaII~{\textsc{H\&K}}. We, therefore, estimate that the emission in CaII~{\textsc{IRT-a}} ranges between 2 and 100\%, this latter value being the maximum ideal.\\

To determine the planetary magnetic field strength ($B_p$) and magnetospheric radius ($r_M$), we considered several SPI models (see Sec.~\ref{Sec:SPI-models}). The only model capable of reproducing the observed power levels is the interconnecting loop model from (\cite{Lanza_2013}). Table~\ref{Tab:Bfields} shows the derived planetary magnetic field and magnetospheric radius of GJ~436~b assuming different stellar magnetic fields and fractions of the total power (see details in Methods).
Under the most conservative assumptions, the planetary magnetic field ranges from 6.3 to 110 G. In comparison, theoretical studies of hot Neptunes orbiting earlier spectral type stars (\cite{Kilmetis_2024}) predict $B_p$ values between 0 G and 11 G, which in our case would correspond to a detected fraction greater than 50\%. The magnetospheric radius associated with these magnetic fields ranges from $6\,R_p$ to $21\,R_p$. This value is smaller than Neptune’s magnetosphere ($\sim 26\,R_{\text{Nep}}$), which is expected given the proximity of GJ~436~b to its host star and the higher stellar wind pressure it receives \cite{Ip_2004, Vidotto_2013}. \\

These constraints on the magnetic field and magnetospheric radius offer valuable insights into the planet’s atmospheric mass loss. Furthermore, the repeated detections of magnetic SPI signals throughout the stellar cycle open new opportunities for future studies and measurements of exoplanetary magnetic fields.



\begin{table}[]
\caption{\textbf{Planetary magnetic field and magnetospheric radius of GJ~436~b}. Magnetic field strengths and magnetospheric radii are derived from the interconnecting loop model equation (Eq.~\ref{p_il}) for different fractions of the detected energy and wind models. $B_p$ is given in Gauss and $r_M$ in planetary radii. Different values are obtained given the assumptions on the stellar magnetic field and are obtained from \cite{Vidotto_2023} and \cite{Lanza_2018} as described in the Methods (Sect.~\ref{subsub:Bp} and \ref{subsub:Rm}, respectively).}\label{Tab:Bfields}



    \begin{tabular}{l|c c|c c|c c}
        \midrule
        \multicolumn{1}{l}{\multirow{2}{*}{\textbf{Models}}} & \multicolumn{2}{c}{\textbf{2 \%}} & \multicolumn{2}{c}{\textbf{25 \%}} & \multicolumn{2}{c}{\textbf{100 \%}}\\
        \multicolumn{1}{c}{} & {$B_p$} (G) & {$r_M$} ($R_p$) & {$B_p$} (G) & {$r_M$} ($R_p$) & {$B_p$} (G) & {$r_M$} ($R_p$) \\ \midrule


        \multicolumn{1}{l} {V23-MI} &  110 & 20 & 24 & 12 & 10.3 & 9   \\
        \multicolumn{1}{l} {V23-MII} & 95 & 13 & 21 & 8 & 9 & 6  \\
        \multicolumn{1}{l} {L18} & 68 & 21 & 14.7 & 13 & 6.3 & 9.6   \\ \hline

    \end{tabular}
\end{table}




\begin{figure}
    \centering
    \includegraphics[width=\columnwidth]{Figure_1.pdf}
    \caption{\textbf{Periodograms of different calcium chromospheric indicators where the SPI signals were detected}. (\textbf{Top}) GLS periodogram of the $I_{CaII}$ index from H-08 epoch. (\textbf{Middle}) Same as Top panel but for pEW'(CaII~{\textsc{IRT-a}}) index from C-16 epoch. (\textbf{Bottom}) Same as Middle panel but for C-24 epoch. In all panels, three vertical lines mark the periods of interest: a dot-dashed line at $P_{\text{SPI-}}=2.81 \,\,\text{d}$, a solid line at $P_{\rm orb}$, and a dotted line at $P_{\text{SPI+}}=2.46 \,\,\text{d}$. A dashed horizontal line shows the 1\% FAP level calculated as described in the Methods. Below each periodogram, the activity indicator time series of the entire epoch is shown. The gray shaded areas correspond to the data used to calculate the periodogram above them, known as sub-epochs. Each panel is color-coded, with blue representing HARPS data and orange representing CARMENES data.}
    \label{Fig_1}
\end{figure}

\begin{figure}
    \centering
    \includegraphics[width=\columnwidth]{Figure_3.pdf}
    \caption{\textbf{Real and modeled observations of C-24 at different phases.} \textbf{(Top):} Generalized Lomb-Scargle (GLS) periodogram of C-24 pEW'(CaII~{\textsc{IRT-a}}) index. \textbf{(Middle):} GLS periodogram of the synthetic data, with the phase of the injected signals adjusted to recreate the observed peaks in the periodogram as described in Section \ref{Modeling-sec}. The left panel shows the phase configuration that recreates the 2.81\,d peak from C-16 and H-08, the middle panel recreates the 2.46\,d peak from C-24, and the right panel displays a phase configuration where several peaks appear. In each panel, three vertical lines mark the periods of interest: a dot-dashed line at $P_{\text{SPI-}}=2.81 \,\,\text{d}$, a solid line at $P_{orb}$, and a dotted line at $P_{\text{SPI+}}=2.46 \,\,\text{d}$. \textbf{(Bottom):} The orange line represents the modeled signal from which the datapoints at the same timestamps as the ones from C-24 are taken. Then, we added an additional Gaussian-distributed noise with a standard deviation of $1\sigma$. The error bars are generated from a Gaussian distribution derived from C-24's original error bars.}
    \label{Fig_2}
\end{figure}



\begin{figure}
    \centering
    \includegraphics[width=\columnwidth]{Figure_2.pdf}
    \caption{\textbf{(Top)}: GJ 436 chromospheric activity indicators from different instruments. HARPS (blue) $I_{CaII}$ index is plotted in the left y-axis. CARMENES (orange dots for C-16 and triangles for C-24)  pEW'(CaII~{\textsc{IRT-a}}) index is plotted together in the right y-axis. \textbf{(Bottom)}: Photometry from different instruments showing the long term photometric magnetic cycle. APT $\Delta (b+y)/2)$ photometry (gray) and AIT (red) after correcting the offset of the difference calculated in the four contemporary epochs and plotted together. The seasonal mean values are plotted for each epoch as large filled circles. The black and dashed line connect the seasonal means of the datasets that show a modulation. The gray shaded regions show the epochs where SPI has been detected (H-08, C-16 and C-24).}
    \label{Fig_3}
\end{figure}

\section*{Acknowledgments}

We thank Giovanni M. Mirouh for the conversations and the ideas that led to determine the role of the phase in the geometrical model.


\textbf{Funding:} We acknowledge the support trough the PRE2021-099654 and PRE2020-095421 grants (D.R. and L.P.), funded by MCIU/AEI/10.13039/501100011033. We acknowledge the support trough the Spanish National grant PID2022–137241NB–C42 (D.R., P.A.)

Support for this work was provided by NASA through the NASA Hubble Fellowship grant HST-HF2-51559.001-A awarded by the Space Telescope Science Institute, which is operated by the Association of Universities for Research in Astronomy, Inc., for NASA, under contract NAS5-26555 (R.L.).

We acknowledge long-term support from NASA, NSF, and Tennessee State University (G.H.).

This research was supported by the ISRAEL SCIENCE FOUNDATION (grant No.\ 1404/22) and the Israel Ministry of Science and Technology (grant No.\ 3-18143) (A.B. and S.Z.).

We acknowledge the support by the European Research Council (ERC) under the European Union’s Horizon 2020 research and innovation programme (ERC Starting Grant "IMAGINE" No. 948582). We also acknowledge the support from ``Mar\'ia de Maeztu'' award to the Institut de Ciències de l'Espai (CEX2020-001058-M) (S.K. and D.V.). Part of the work has been carried out within the framework of the doctoral program in Physics of the Universitat Autònoma de Barcelona (S.K.).

\textbf{Author contributions:} D.R. and P.A. conceived the idea. D.R. led the analysis of the data. D.R., P.A and R.L interpreted the data. D.R., R.L. and P.A. wrote the manuscript. P.S. obtained CARMENES activity indicators and helped to interpret the data. N.L. developed the geometrical model. A.H. contributed to the analysis of the significance of the signal. G.H. provided and helped in the analysis of APT and AIT data. A.B. and S.Z. analyzed the spectra with PDC, \texttt{USuRPER} and \texttt{SPARTA}.
The remaining authors including the already mentioned provided comments and reviewed the manuscript.
\textbf{Competing interests:} The authors declare no competing interests.
\textbf{Data and materials availability:} The HARPS spectroscopic observations are available at: \url{http://archive.eso.org/eso/eso_archive_main.html}. The CARMENES 2016 spectroscopic observations are available at: \url{http://carmenes.cab.inta-csic.es/gto/jsp/dr1Public.jsp.}


\clearpage
\newpage

\section*{Methods}\label{sec2}
\subsection{Observations and data reduction}
\subsubsection{HARPS spectroscopy} \label{HARPS}
There are 192 observations of GJ 436 taken with the HARPS instrument (\cite{Mayor_2003}) at the 3.6m ESO telescope in La Silla, Chile, between January 2006 and March 2020. The data is analyzed in \cite{Kumar_2022} and we follow the same methodology for reducing the data and obtaining the activity indicators. The 2008 epoch, where SPI was detected, consists of 47 observations. As described in \cite{Kumar_2022}, the HARPS-reduced spectra were not continuum-normalized and were used to derive magnetic activity indices for CaII H\&K, HeID3, NaID, $H\alpha$ and CaI.\\

We use the \texttt{ACTIN2} package (\cite{da_Silva_2011} ,\cite{da_Silva2018}, \cite{da_Silva_2022}) to calculate the CaII H\&K index as:
\begin{equation}
    I_{\rm CaII} = \frac{F_{\rm CaII\,K} + F_{\rm CaII\,H}}{F_1 + F_2}
\end{equation}\\
where $F_{CaIIK}$ and $F_{CaIIH}$ represent the mean fluxes integrated in triangular windows around the K and H line and $F_1$ and $F_2$ are the mean fluxes in square windows of the pseudo-continuum. Further details on the calculation of this and other indices are provided in Section 3 of \cite{Kumar_2022}.

\subsubsection{CARMENES spectroscopy} \label{CARMENES}
GJ 436 was observed with the CARMENES instrument (\cite{Quirrenbach_2013},\cite{Quirrenbach_2014}) at the 3.5\,m telescope of the Calar Alto Observatory in Almeria, Spain. There are 444 observations taken between January 2016 and April 2021, of which only 112 (with an average cadence of one observation every two days) were used for this study. The remaining observations, primarily high-cadence, intranight observations were obtained for atmospheric studies. We revisited GJ 436 in 2024 (program ID: 24A-3.5-013) to confirm the tentative detection of SPI seen in HARPS-2008 and CARMENES-2016 data, as the star was in the same phase of its magnetic cycle. This additional program collected 51 observations from January 1 to May 16, 2024, with an average cadence of one observation every three days. We derived pseudo-equivalent widths (pEW') of the He D3, Na D, H$\alpha$, and Ca IRT lines (being the \textsc{IRT-a} the bluest of the three) using a spectral subtraction technique and indices of selected TiO and VO absorption bands from the reduced spectra as described in \cite{Schofer_2019}.

\subsubsection{HIRES spectroscopy} \label{HIRES}
There are 274 additional observations of the GJ 436 system between 2004 and 2019 taken with the HIRES instrument (\cite{Vogt_1994}) installed on the Keck I telescope at Mauna Kea, Hawai`i. Two independent studies have used those observations to calculate the S-index (\cite{Rosenthal_2021}) and the CaII~{\textsc{H\&K}} (\cite{Perdelwitz_2021}) of GJ 436. Both datasets were considered in our analysis of stellar activity indicators. However, the low signal-to-noise ratio of most of the observations made them unsuitable for detecting any signals. Therefore, the HIRES observations were ultimately excluded from this study.

\subsubsection{APT and AIT photometry}
We have acquired 2040 photometric observations of GJ~436 during 15 observing seasons between 2003 and 2018 with the Tennessee State University (TSU) T12 0.80\,m Automatic Photoelectric Telescope (APT) at Fairborn Observatory in southern Arizona.  The first 14 observing seasons were discussed in \cite{Lothringer_2018}.  The T12 APT is equipped with a two-channel photometer that uses two EMI 9124QB bi-alkali photomultiplier tubes to measure stellar brightness simultaneously in the Str\"omgren $b$ and $y$ passbands.  A complete description of the TSU T8 APT, which is identical in construction and operation to T12, can be found in \cite{Henry_1999}.\\

The observations of GJ~436 (star d) were made differentially with respect to three nearby (on the sky) comparison stars:  HD~102555 (star a), HD~103676 (star b), and HD 99518 (star c).  We measure the difference in brightness between our program star GJ~436 (d) and the comparison stars a, b, and c and create differential magnitudes in the following six combinations:  d-a, d-b, d-c, c-a, c-b, and b-a.  Intercomparison of these six light curves shows that comparison stars a and b both appear to be constant, so we present our results as differential magnitudes: star~d minus the mean brightness of stars a and b (d-ab).  To improve the photometric precision of the individual nightly observations, we combined the differential $b$ and $y$ magnitudes into a single $(b+y)/2$ ``passband''.  The observations are all corrected for differential atmospheric extinction and transformed to the standard Str\"omgren $b$ and $y$ passbands using nightly observations of a set of standard stars distributed throughout the sky.\\

Our T12 observations are plotted in the top panel of Fig. \ref{Fig:APT-Greg} as small filled circles.  The mean of the nightly observations is plotted as the dashed line.  The differential magnitudes of the two comparison stars are plotted in the bottom panel of Fig. \ref{Fig:APT-Greg}. The standard deviation of the 2043 nightly comparison star observations from their mean is 2.43 mmag in the $by$ ``passband'', providing an estimate of the precision of our individual differential magnitudes of GJ~436.  Comparison of the two panels of Fig. \ref{Fig:APT-Greg} shows that GJ~436 exhibits low-amplitude night-to-night and year-to-year brightness variability.\\

We also acquired 607 additional observations in the Cousins $R$ band with the TSU Celestron 14-inch (C14) Automated Imaging Telescope (AIT) covering the eight observing seasons 2013 to 2017 and 2020 to 2024. The AIT uses an SBIG STL-1001E CCD camera with a Kodak KAF-1001E detector and gives a 21 x 21 arcmin field of view.  Each nightly observation consisted of 4–10 consecutive exposures centered on GJ~436.  The individual nightly frames were co-added and reduced to differential magnitudes as GJ~436 minus the mean brightness of four constant comparison stars in the same field of view.  Each nightly observation has been corrected for bias, flat-fielding, pier-side offset, and differential atmospheric extinction.\\

The observations are plotted in Fig. \ref{Fig:C14-Greg}. The mean of the 8 seasons is represented by the dashed line in the plot. The first four observing seasons overlap with the T12 observations while the last four seasons extend our observations by several years. The C14 observations show the same low-amplitude nightly and yearly brightness variations as seen in the T12 data.\\

\begin{figure}
    \centering
    \includegraphics[width=\columnwidth]{T12-Greg.pdf}
    \caption{\textbf{T12 APT Photometry of GJ 436}. \textit{Top:} nightly Str\"omgren $(b+y)/2$ band photometry of GJ~436 (small filled circles). They scatter about their mean (dashed line) with a standard deviation of 5.62~mmag. \textit{Bottom:} differential magnitudes of the comparison stars at the same scale.  The comparison star observations scatter about their mean with a standard deviation of 2.43 mmag.  A comparison of the two panels indicates low-amplitude variability in GJ 436 on both night-to-night and year-to-year timescales.}
    \label{Fig:APT-Greg}
\end{figure}

\begin{figure}
    \centering
    \includegraphics[width=\columnwidth]{C14-Greg.pdf}
    \caption{\textbf{C14 AIT Photometry of GJ 436}. Cousins $R$ band photometry during the observing seasons 2013-14 to 2016-17 and 2020-21 to 2023-24.  Like the T12 observations, they show low-amplitude variability on both night-to-night and year-to-year timescales.}
    \label{Fig:C14-Greg}
\end{figure}




\subsection{Periodograms, tools, and procedures}
\subsubsection{Splitting datasets}\label{subsub:GLSchunk}

The search for SPI is an active research topic in the exoplanet field, particularly in the optical domain (\cite{Shkolnik_2003},\cite{Shkolnik_2004},\cite{Shkolnik_2008},\cite{Fares_2012},\cite{Pillitteri_2011},\cite{Pillitteri_2014},\cite{Pillitteri_2015},\cite{Cauley_2017},\cite{Cauley_2018},\cite{Klein_2022},\cite{Kumar_2022},\cite{Ilin_2023}, \cite{Loyd_2023}). One critical issue arises from the uncertainty regarding when and for how long an SPI signal might manifest. This lack of knowledge complicates the search, as the detectable signal may only be present at certain epochs. A promising approach involves splitting the dataset into smaller chunks, which reduces the risk of diluting any transient signal that could otherwise be lost in the full dataset. This method allows us to isolate moments where the SPI signal is potentially detectable, focusing on smaller windows of time where the interaction might be more evident.\\



We apply this procedure to our spectroscopic datasets and search for periodicities using a GLS periodogram \cite{Zechmeister_2009}. The results are shown in Figs.~\ref{HARPS_2008full_All} through \ref{CARMENES_2024_All}. For HARPS, we divide the dataset in multiple chunks corresponding to the visibility of the target (January-June). Out of the six epochs and from all the activity indicators, we only detected a signal compatible with SPI at $P_{\text{SPI-}} = 2.81\,\text{d}$ in the CaII~{\textsc{H\&K}} from 2008 (H-08). The false alarm probability (FAP) of this signal is near 1\% which is still moderately significant. For CARMENES, all the data was kept in two chunks corresponding to the 2016 (C-16) and 2024 (C-24) observing seasons. In the C-16 chunk, just like in the H-08 data, we find a single moderately significant peak at $P_{\text{SPI-}} = 2.81\,\text{d}$ in the CaII~{\textsc{IRT-a}} line with a FAP near 10\%. The C-24 observations have a much better time sampling and we find a signal at $P_{\text{SPI+}} = 2.46\,\text{d}$ with a FAP below 1\% in the CaII~{\textsc{IRT-a}} line.



\subsubsection{Aliasing}\label{Sec:AliasProblem}

\begin{figure}
    \centering
    \includegraphics[width=0.99\linewidth]{AF_GJ436_H08.pdf}
    \caption{\textbf{Aliasing analysis with \texttt{AliasFinder}.} \textit{Top:} The original data periodogram (red) plotted against the median GLS periodogram (dark) from simulations where a signal of 2.81\,d was injected (vertical blue dashed line). The other two columns show the region around the 1-day aliases: 1.55\,d and 0.73\,d. \textit{Middle and bottom:} The same as the top row but injecting a signal of 1.55\,d and 0.73\,d, respectively. The shaded grey areas show the interquartile range and the range of 90 and 99\% of the simulated periodograms. The clocks show the angular mean of the phase of each peak with the standard deviation.}
    \label{fig:AliasFinder_H08}
\end{figure}

\begin{figure}
    \centering
    \includegraphics[width=0.99\linewidth]{AF_GJ436_C16.pdf}
    \caption{\textbf{Same as Fig. \ref{fig:AliasFinder_H08} but for the C-16 dataset.}}
    \label{fig:AliasFinder_C16}
\end{figure}

\begin{figure}
    \centering
    \includegraphics[width=0.99\linewidth]{AF_GJ436_C24.pdf}
    \caption{\textbf{Same as Fig. \ref{fig:AliasFinder_C16} but for the C-24 dataset, injecting a signal of 2.46\,d  1.68\,d and 0.71\,d }}
    \label{fig:AliasFinder_C24}
\end{figure}
The periodograms of the calcium indicators from all three analyzed epochs reveal what we interpret as the primary SPI signals (H-08: 2.81\,d, C-16: 2.81\,d, C-24: 2.46\,d) and their potential $\pm$~1-day aliases. In H-08, the alias is slightly stronger than the primary signal, while in C-16 and C-24, the primary signal dominates. To validate our interpretation of the primary signals and aliases, we use the code \texttt{AliasFinder} \cite{Stock_2020}.
This method is based on the work from \cite{DawsonFabrycky2010ApJ...722..937D} and has been used in \cite{Trifon_2018} to determine the true signals in radial velocity attributed to the planetary companion GJ~436~b (2.64\,d vs. its 1-day alias at 1.6\,d). The relatively weak power of the peaks in all the three epochs analyzed did not yield any significant results using \texttt{AliasFinder} (Figs.~\ref{fig:AliasFinder_H08},\ref{fig:AliasFinder_C16},\ref{fig:AliasFinder_C24}).



\subsubsection{Rolling periodograms}\label{Sec:Rolling}
The previous approach was useful in identifying signals of potential interest, but their statistical evidence is not enough to be confirmed as bona fide periods present in the data. To overcome the limitations of standard periodograms to find transient signals, we use rolling periodograms \cite{Herbort_2018,Schofer_2021}. Rolling periodograms are periodograms calculated from overlapping data subsets. For our data, we keep the chunks from each instrument as in the previous section and then divide those into smaller sub-epochs of a desired length, computing the GLS periodogram for each of them. If the epoch has $N$ data points and we define each sub-epoch to have $m$ points, where i is the index of the observation's position in the subset, the first sub-epoch will be $[i_0, i_m]$. The next sub-epoch will be $[i_1, i_{m+1}]$, followed by $[i_2, i_{m+2}]$, and so on, until reaching $[i_{N-m}, i_{N}]$. This allows us to capture variations in the signal across different parts of the split dataset. \\

For HARPS, out of the 6 chunks two were discarded due to having too few observations ($<10$) to carry out this analysis. In the H-08 data, our rolling periodogram analysis finds the signal at $P_{\text{SPI}-}$ to appear partway through the dataset, peak between February 26 and March 11, and then disappear (Fig.~\ref{fig:Rolling_periodogram_HARPS}). In the C-16 data, the signal has a similar behavior, with the $P_{\text{SPI}-}$ appearing towards the end of the observation window  but with less intensity. For the C-24 dataset, due to its shorter time span, we do not find differences between the periodogram of the full dataset or the rolling periodogram. We analyze the significance of all these signals in the next Section.



\subsubsection{Significance of the signals}\label{subsub:FAP}

The FAP levels reported above consider the full dataset, but as demonstrated with the rolling periodograms, the signals of interest are not present at all times. Moreover, in our analysis we are not blindly searching for periodicities in the full frequency space but rather looking into a certain frequency range. To estimate the FAP levels of finding a signal at the expected periods, we perform two independent bootstrap analyses (\cite{Murdoch_1993}) in the close range of periods around the 2.81-day signal in the H-08 and C-16 data, and 2.64-day signal in the C-24 data. We compute the GLS periodogram of 10$^4$ time series obtained by keeping the time stamps of the data and, for each of them, selecting a random flux value of the data with replacement.\\

As the peaks and the power in the GLS periodogram are dependent on the range of periods it is computed, we set as a reference the FAP of the 2.81-day signal found in H-08 and C-16 and the 2.46-day signal from C-24 in Table \ref{FAP-table} that we name FAP$_{\rm GLS}$. The [0.7\,d - 10\,d] range is the one used in the previous section, while the narrower [1.4\,d - 3\,d] was chosen to isolate the 2.81\,d/2.46\,d periods. Then, for each of the $10^4$ simulated series in those two ranges, we compute the GLS periodogram and check how many times a peak higher than the 2.81\,d/2.46\,d peak in the original time series and period range appeared, giving us the FAP$_{\rm boot}$ value. We perform this search for two period ranges for each of the previously defined period ranges and instruments, corresponding to 3 and 1.5 times the FWHM of the signal of interest. For H-08, this corresponds to [2.49\,d - 3.21\,d] and [2.64\,d - 3\,d], respectively. For C-16 and C-24 data, these intervals correspond to [2.7\,d - 2.9\,d] and [2.75\,d - 2.85\,d] and [2.27\,d - 2.71\,d] and [2.34\,d - 2.61\,d], respectively.

The FAP$_{\rm boot}$ from Table \ref{FAP-table} represents the probability that a random peak appears in those period intervals as a result of the noise in the data. The joint probability of the signal appearing in those periods, in three different epochs, and with that significance, is the multiplication of the individual FAP$_{\rm boot}$ probabilities. Conservatively, we select the largest FAP$_{\rm boot}$ values for each epoch (H-08: 2.37\%, C-16: 2.3\%, and C-24: 0.16\%) to obtain a combined FAP of $8.7 \cdot 10^{-6}$. Therefore, the probability that a source of noise would produce such a signal in three different epochs would be less than one in 10$^5$ realizations, making the case for this signal to be statistically significant despite its low significance in each individual chunk and instrument.

\begin{table}
\caption{\textbf{Results from the signal significance analysis in Sect.~\ref{subsub:FAP}.} False alarm probabilities from the GLS (FAP$_{\rm GLS}$) and the bootstrap (FAP$_{\rm boot}$).}\label{<table-label>}
\begin{tabular}{c c c c c c}
\toprule
Instrument & Sub-epoch & Interval [d] & FAP$_{\rm GLS}$ [\%] & Interval [FWHM] &  FAP$_{\rm boot}$ [\%] \\
\midrule
HARPS & 2008 & 0.7 -10 & 3.4 & 3  &  2.37 \\
 &  &  &  & 1  & 1.37 \\
 &  & 1.4 - 3 & 1 & 3 & 2.13 \\
  &  &  &  & 1 & 1.32 \\
CARMENES & 2016 & 0.7 -10 & 35 & 3  & 2.3 \\
 &  &  &  & 1  &  0.86 \\
 &  & 1.4 - 3 & 11.77 & 3  & 1.7 \\
 &  &  &  & 1  & 0.95 \\
CARMENES & 2024 & 0.7 -10 & 0.27 & 3  &  0.13 \\
 &  &  &  & 1 &  0.11 \\
 &  & 1.4 - 3 & 0.077 & 3 & 0.16 \\
 &  &  &  & 1 & 0.1 \\
\botrule
\end{tabular}

\label{FAP-table}
\end{table}




\subsubsection{Other periodograms}
\paragraph{$\ell$-1 periodogram}
The $\ell$-1 periodogram (\cite{Hara_2016}) is designed to detect periodic signals in unevenly sampled time series, similar to the Lomb-Scargle (\cite{Scargle_1982}) periodogram, but less prone to aliasing. Primarily developed for exoplanet detection in radial velocity data, it can also be applied to other fields requiring periodicity detection. We use this periodogram to analyze the split data and find identical results to those obtained from the GLS in Sect.~\ref{subsub:GLSchunk}, thus opting for not reporting it here for the sake of brevity.

\paragraph{PDC, \texttt{USuRPER} and \texttt{SPARTA} periodograms}
The Phase Distance periodogram (PDC) (\cite{Shay_2017}) is a model-independent tool for periodicity detection in astronomical time series, designed to detect general periodicities in unevenly sampled data. \cite{Binnenfeld_2020} extended it via \texttt{USuRPER}, a method for detecting periodic variation in time series of astronomical spectra. The novelty of \texttt{USuRPER} lies in its ability to detect periodic variations in the overall spectra directly, without relying on specific scalar quantities like radial velocities or activity indicators. This method is ideal for identifying and analyzing signals associated with stellar activity and pulsations. \cite{Binnenfeld_2020} further enhanced this approach, enabling the separation of spectral shape variations from bulk Doppler shifts related to orbital motion.\\

We used SPectroscopic vARiabiliTy Analysis (\texttt{SPARTA}) (\cite{Shahaf_2020}), a suite of tools that includes \texttt{USuRPER}, to apply these periodograms to all the three epochs H-08, C-16 and C-24 and found no signal at the SPI$\pm$ periods. The results from this analysis showed periodic signals at 2.24 and 1.79\,d  corresponding to non-Doppler spectral variations. These signals are not alias of each other and are probably not related with our SPI signals.\\

\begin{comment}
\begin{figure}
    \centering
    \includegraphics[width=1\linewidth]{Fig_Avraham.pdf}
    \caption{\textbf{PDC periodograms for the HARPS spectra of GJ 436.} \textbf{Top panel}: The Shape periodogram displays two prominent peaks at 2.24 and 1.79 days (indicated by vertical blue dashed lines), strictly corresponding to non-Doppler-induced variations in the line profiles. \textbf{Bottom panel}: The Shift Periodogram shows no significant peaks within the trial frequency range, further supporting a stellar activity origin for the detected signals. The dotted horizontal line in both panels represents a false alarm probability (FAP) threshold of $10^{-3}$.}
    \label{fig:PDC}
\end{figure}
\end{comment}


\subsection{Modeling and interpretation of the signal} \label{Modeling-sec}

\begin{figure}
    \centering
    \includegraphics[width=0.8\linewidth]{Coordinates.pdf}
    \caption{\textbf{Coordinate system and reference frame used in the model.} The sphere represents the stellar surface, with the XYZ coordinate system centered on GJ 436 (O), with the Z-axis towards the rotation axis. The stellar equator is contained in the XY plane, and the blue dot (A) represents the position of a stellar active region described in the model. The dashed circle represents the projection of GJ~436~b orbit on the surface of the star. The orbit is contained in the xy plane and has an inclination I. The ascending node (N) is represented as a dark spot and the x-axis points towards it. A red dot (P) represents the position of an induced chromospheric hot spot that follows the planetary orbit. AP is the distance between the chromospheric hot spot induced by the planet (red) and the one intrinsic to the stellar activity (blue).}
    \label{fig:Coordinate}
\end{figure}

In this section, we propose a model to explain why the signal appears at two distinct, yet related, periods/frequencies that differ from the planetary orbital period itself. This geometrical model is based on the hypothesis that the intensity of the signal depends on the distance of the signals coming from two active regions on the chromosphere of the star. One of these active regions has a stellar origin and thus rotates with the star (point A), while the other is induced by SPI and moves according to the planet’s orbit (point P). The combined observations of those two separated and independent regions would produce the splitting from $P_{\text{orb}}$ to $P_{\text{SPI-}}$ and $P_{\text{SPI+}}$.\\

Let us consider a reference frame with its origin $O$ at the center of the star and the $XY$ plane coincident with the equatorial plane of the star. The orbit of GJ~436~b is highly inclined and moderately eccentric. For simplicity, we consider only the effect of the obliquity. The obliquity of a planetary orbit $I$ is the angle between the equatorial plane of the star and the plane of the orbit, which corresponds to  $I = 103^{\circ} \pm 13^{\circ}$ for GJ~436~b. This value was measured by \cite{Bourrier_2018} and is the true 3D obliquity, not the obliquity projected on the plane of the sky (typically referred as $\lambda$). We assume that the $X$ axis is directed along the line of nodes of the planetary orbit pointing towards the ascending node. Considering the coordinate transformation shown in Fig.~\ref{fig:Coordinate}, and assuming the orbit to be circular with an orbital radius $a$, the coordinates of the planet $P(X_{\rm p}, Y_{\rm p}, Z_{\rm p})$ are

\begin{equation}
    \left\{
\begin{array}{lll}
X_p &= a \cos (f+\phi_1), \\
Y_p &= a \sin (f+\phi_1) \cos I, \\
Z_p &= a \sin (f+\phi_1) \sin I
\end{array}
\right.
\end{equation}
where $f = n t$ is the orbital phase of the planet along its orbit with $n \equiv 2\pi/P_{\rm orb}$ being its mean angular frequency, $P_{\rm orb}$ the orbital period, and $t$ the time measured from the passage through the ascending node.\\

Assuming the active region $A(X,Y,Z)$ is located on the equator of the star, its coordinates are

\begin{equation}
    \left\{
\begin{array}{lll}
X &= R \cos (\lambda +\phi_2), \\
Y &= R \sin (\lambda +\phi_2), \\
Z & = 0
\end{array}
\right.
\end{equation}
where $R$ is  the radius of the star, $\lambda = \lambda_{0} + \Omega t$ with $\lambda_{0}$ the longitude of the point $A$ at time $t=0$, $\Omega = 2\pi/P_{\rm rot}$ the stellar spin angular velocity, and $P_{\rm rot}$ the stellar rotation period. One of the two phases ($\phi_1$ or $\phi_2$) can always be made zero by an appropriate choice of the origin of the coordinate system; in this case, we set $\phi_2 = 0$.\\

Under our hypothesis, the intensity of the SPI signal depends on the distance $d = \overline{AP}$. Substituting the coordinates of the points $A$ and $P$ into the distance formula, we find after a few trigonometric and algebraic manipulations:

\begin{equation}
    \begin{array}{l}
 \overline{AP}^{2} = (X-X_{\rm p})^2 + (Y-Y_{\rm p})^{2} + Z_{p}^{2} = \\
=  R^{2} + a^{2} -2aR \left( \cos \lambda \cos (f+\phi_1) + \sin \lambda \sin (f+\phi_1) \cos I \right) = \\
= R^{2} + a^{2} - 2aR \cos (\lambda-f-\phi_1) +2 aR (1 - \cos I) \sin \lambda \sin (f+\phi_1).
    \end{array}
\end{equation}\\

The difference $\lambda - f = (\Omega -n) t + \lambda_{0}$, while the product $\sin \lambda \sin f =  \left[ \cos(\lambda-f) - \cos(\lambda+ f) \right]/2$ by applying Werner's formulae. In conclusion, we can write:

\begin{equation}
\begin{array}{l}
    \overline{AP}^{2} = R^{2} + a^{2} - 2 aR \cos \left[ (\Omega-n) t + \lambda_{0 }-\phi_1 \right] + \\
    aR (1 -\cos I) \left\{ \cos \left[ (\Omega-n)t + \lambda_{0}-\phi_1 \right] - \cos \left[ (\Omega+n) t +\lambda_{0}+\phi_1 \right] \right\}.
\end{array}
    \label{dist2}
\end{equation}\\

Finally, if we make $R=a$ the expression would describe the motion of a point on the star that is following the motion of the planet along its orbit. This is called the sub-planetary point on the star. Then, Eq.~\ref{dist2} shows that the distance $\overline{AP}$ oscillates with the synodic frequency $\Omega -n$ when the obliquity of the orbit is zero because $I = 0$ implies $\cos I =1$. On the other hand, for an oblique orbit as in the case of GJ~436~b, $I \not= 0$, and the other frequency $\Omega + n$ appears. Therefore, the appearance of the frequency $\Omega +n$ in the SPI periodicity of GJ~436 should be related to the obliquity of its orbit. This explains why that frequency was not previously observed in other systems where the orbit is aligned with the equatorial plane of the host star.\\

The orbit of GJ~436~b is eccentric with an eccentricity $e = 0.152$ (\cite{Vidotto_2023}). Introducing the eccentricity in the above model is possible, but it complicates the mathematical development because both the orbital radius $r(t) \equiv OP$ and the longitude of the planet along its orbit $f$ become complex functions of the time. Nevertheless, given that those functions are periodic in $P_{\rm orb}$, they can be developed in Fourier series and we can repeat the above analysis for each of the coefficients of the series corresponding to the harmonics of the orbital frequency $n$. The dominant terms in the Fourier series of $r$ and $f$ are those with the fundamental orbital frequency $n$ because the eccentricity of GJ~436~b is not large, so the higher harmonics have smaller and smaller amplitudes as their order increases. In conclusion, our model assuming a circular orbit can be regarded as an approximation of the first order.

\subsubsection{Alternating occurrence of the signals at $\Omega - n$ and $\Omega + n$}

We find that the reason why the signal sometimes appears at $\Omega - n$ and at other times at $\Omega + n$ is apparently due to the phases of the two interacting components, which are ultimately two regions (one on the star and the other caused by the planet) that interact and have periodicities of $P_{\text{rot}}$ and $P_{\text{orb}}$, respectively.\\

To explore the behavior of Eq.~\ref{dist2}, we generate synthetic data adding noise drawn from a Gaussian distribution with a standard deviation of $1\sigma$. The error bars were generated from a separate Gaussian distribution with the same standard deviation as the original error bars in the C-24 data. The time variable was selected using the same timestamps as in the C-24 observations.\\

By keeping the phase $\phi_1$ as a free parameter, we observe that varying its value and generating the periodogram of the resulting data reproduces observations where only one of the two peaks appears at a time (Fig.~\ref{Fig_2}). Thus, regardless of phase, the signal at $P_{\rm orb}=2.64$\,d that corresponds to the planet's period does not appear in the periodogram. Instead, this signal is split into the other two, $P_{\text{SPI}-}$ and $P_{\text{SPI}+}$. The other period of interest, $P_{\rm rot}$, could appear in the interaction's resulting periodogram as seen in our synthetic data (Fig.\ref{Fig_2}). However, we do not find this signal in the calcium indicators, but in the H$\alpha$ line.\\

The reason behind the change in the phase between epochs might be directly related with the nature of the two regions themselves. Although the planet moves at a repetitive and predictable path, the hot spot induced by its interaction with the stellar magnetic field can lag or overtake the subplanetary point. The magnetic fields lines involved in the SPI between both bodies can have different geometries and shapes, thus a direct, straight and constant connection between the planet and the star is not expected. This difference was already seen in the first SPI tentative detections (\cite{Shkolnik_2003}). On the side of the star, active regions are known to change in shape, intensity and even latitude trough the activity cycle (Spörer's law). Therefore, it is unlikely that a constant phase between both regions is occurring.


\subsection{Stellar properties}

\subsubsection{Activity cycle} \label{Activity_cycle}
Understanding the magnetic cycle of GJ 436 is crucial, as the stellar magnetic field serves as the primary mediator of star-planet interactions (SPI). Variations in magnetic field strength, stellar wind, and the surrounding stellar medium throughout the activity cycle could significantly influence the detection and characterization of SPI. Consequently, referencing SPI observations within the context of the activity cycle phase may enhance detection likelihood and provide insights into why SPI signals are not consistently detectable.\\

\cite{Lothringer_2018} reported the first indication of an activity cycle in GJ 436 using photometry from APT and revealing a modulation of 7.4 years. Later, studies using the same photometric data determined a slightly longer period of $7.75 \pm 0.10$ years \cite{Loyd_2023}. From HARPS spectroscopic observations, which we also used to study SPI signals, \cite{Kumar_2022} suggested a period of approximately 6.8 years.
However, the addition of a new APT epoch (2017-18) shows a possible change in the periodicity of the cycle. This data shows a decrease in GJ 436's magnitude which could imply an increase in the activity and spot coverage of the star. Contemporary and new observations were taken in another passband (Cousins $R$) by AIT. Even though the APT and AIT observe in different bands, we combine both datasets by adding an offset to AIT data calculated from the mean values of the contemporary epochs. The combined light curve (Bottom panel in Fig.~\ref{Fig_3}) shows a linear trend toward lower magnitudes. Furthermore, the amplitude of the modulation  seen in APT data is not apparent in the AIT data after 2020. Since we have contemporary monitoring of the star during that time, we can ascribe this change to the star rather than the instrument.\\

If we compare the photometric and chromospheric long-term time series, we found an anti-correlation between photometric variability and chromospheric activity ($\Delta\phi \not\approx 0$) (Fig.~\ref{Fig_3}), in line with previous studies \cite{dos_Santos_2019,Loyd_2023}. This anti-correlation suggests a spot-dominated nature. In spot-dominated stars, photometric brightness is at its minimum during the activity maximum, as a larger portion of the stellar surface is covered by spots, which emit less light in the wavelengths typically observed in photometry. Conversely, activity indicators, such as the CaII~{\textsc{H\&K}} lines, reach their maximum, consistent with our observations of GJ 436 (Fig.~\ref{Fig_3}).


\subsubsection{Rotation period} \label{Sec:Rotation}
There have been numerous estimates of GJ 436's rotation period thanks to the extensive monitoring of this system. Photometric observations with the APT over 14 years revealed a rotation period of $P_{\rm rot} = 44.09 \pm 0.08$\,d (\cite{Bourrier_2018}). A variability in the period was observed across different epochs, ranging between 41.7\,d to 46.6\,d (\cite{Bourrier_2018}). Subsequent independent analyses of these observations confirmed this rotation period (\cite{Lothringer_2018}, \cite{Loyd_2023}). Additionally, results from SuperWASP photometric observations (\cite{Pollaco_2006}) showed a rotation period of $P_{\rm rot} = 44.6 \pm 2$\,d (\cite{DiezAlonso_2019}). Furthermore, spectroscopic activity indicators such as $H\alpha$ and CaII~{\textsc{H\&K}} confirmed the photometric rotation periods, with values ranging from 40 to 50 days (\cite{Bourrier_2018}, \cite{Lothringer_2018}, \cite{Kumar_2022}). \\

In our analysis, we identify a period consistent with the rotation period in the $H\alpha$ from C-24 at $P_{\rm rot} = 40 \pm 2.7$\,d. However, the significance of this peak only reached the $7\%$ false alarm probability (Fig.~\ref{CARMENES_2024_All}). No other periods compatible with the rotation period were found in the activity indices from the 2024 epoch or in any other indicator from C-16 (Fig. \ref{CARMENES_2016_All}).



\subsubsection{CaII~{\textsc{H\&K}} and CaII~{\textsc{IRT-a}} fluxes} \label{Sec: powers}
To estimate the planetary magnetic field from SPI, we must determine the power emitted in the SPI interaction from the spectroscopic observations. There are several ways to obtain the power from the spectra. One approach consists of using PHOENIX model spectra to flux calibrate the line variations where the SPI is detected and then convert them to disk-averaged powers (\cite{Cauley_2019}). For our work, we use the calibrations by (\cite{Labarga_2022}) of the $\chi$ factor (\cite{Walkowicz_2004}) for the line where we detected SPI. This approach is simpler, but limited to having pEW's of the lines. For that reason, we only calculate the power emitted during the SPI interaction from CARMENES observations of the CaII~{\textsc{IRT-a}} line.\\

The $\chi$ factor relates the fraction of emitted luminosity with respect to the bolometric luminosity to the equivalent width of the line in the form $\frac{L_{\text{line}}}{L_{\text{bol}}} = \chi \cdot {\rm pEW'}$. Using the calibrations in (\cite{Labarga_2022}), which require knowing the star's effective temperature ($T_{\rm eff}=3533\pm26$K, \cite{Marfil_2021}), we can derive the $\chi$ factors for the CaII~{\textsc{IRT-a}} to be $\chi_{\text{CaII~{\textsc{IRT-a}}}} = 7.99 \times 10^{-5}$. Then, we take the difference between the smallest and the largest values from the pEW' time series in Fig.~\ref{Fig_1}, as numerically smaller pEW' values indicate more emission or less absorption. We find $\Delta {\rm pEW'} = -0.0434$ and 0.048 for the C-16 and C-24 datasets, respecitvely. \\

From the previous equation, we can now determine $L_{\rm line}$ using Stefan-Boltzmann law for $L_{\rm bol} = 4\pi R_{\star}^2 \sigma T_{\text{eff}}^4$. Given \(R_{\star} = 0.41\,R_{\odot}\) and \(T_{\text{eff}} = 3533\,\text{K}\), the emitted power of SPI in the CaII~{\textsc{IRT-a}} line is $L_{\text{C-16}} = 3.14 \times 10^{19}\,\text{W}$ and $L_{\text{C-24}} = 2.23 \times 10^{19}\,\text{W}$ in C-16 and C-24 observations, respectively.



\subsubsection{Fraction of SPI detected power}\label{Sec:Frac_Power}
One of the most important parameters to determine when estimating the total energy of the SPI is the fraction dissipated in the lines where the emission is detected. Various authors have addressed this problem by comparing SPI emission to that produced by solar flares (\cite{Cauley_2019}) or flares from other stars (\cite{Loyd_2023}). Although comparing the SPI emission to the ones in flares can be a first approximation, the real emission might be more like an auroral-like spot, as seen in the Jupiter-Io system (\cite{Dulk_1965}, \cite{Piddington_1968}). This idea arises from the observation that, despite having multiple spectral lines and activity indicators, we only detect a signal attributable to SPI in the CaII~{\textsc{H\&K}} and CaII~{\textsc{IRT-a}} lines. Therefore, it is likely that SPI emission is not as panchromatic as flare emission. Nevertheless, the estimation using flares serves as a lower limit for the fraction of energy emitted in these lines and we follow this approach in our analysis below.\\

Following the multiwavelength observations of AD Leonis by \cite{Hawley_2003}, we find that, in the best case scenario, the emission in the CaII~{\textsc{IRT-a}} line constitutes about 2\% of the total emission during the decay phase, with the CaII~{\textsc{H\&K}} emission being half that. We compare the emission with the decay phase of the flare rather than the impulsive one, as based on the Jupiter-Io system, this type of interaction is expected to be steady and longer-lasting rather than short-lived and eruptive. In the case of AD Leonis, most of the energy radiated in optical lines is concentrated in H$\alpha$. Since we do not detect any signal in this line in our observations, we infer that the energy distribution of the SPI is less broad than in flares, where a larger fraction may be dissipated in the CaII~{\textsc{H\&K}} and CaII~{\textsc{IRT-a}} lines. For a more detailed estimation of the fraction of the total power emitted by SPI in a given line, we require multiwavelength observations during a detection. In our subsequent analyses, we compute the planetary magnetic field and the magnetospheric radius assuming different fractions ranging from 2 to 100\% given the uncertainties described above.



\subsection{Planetary magnetic field estimation}

\subsubsection{Energetics of different magnetic SPI models} \label{Sec:SPI-models}

Magnetic SPI offers the opportunity to estimate planetary magnetic field strengths by comparing the observed energy outputs from the interaction with theoretical predictions from various models. A handful of analytical models have been proposed to quantify the power released through magnetic SPI mechanisms. Here, we follow similar steps as the ones in \cite{Cauley_2019} to obtain $B_p$ from different models, including Alfvén waves, magnetic reconnection, changes in the helicity, and an interconnecting loop between both bodies. All of these mechanisms give different frameworks to explain the interaction and account for different magnitudes of the total power released on it. In this section, we summarize each model and compute the power available using the parameters of the GJ 436 system.

\paragraph{The Alfvén Wing Model}
This model, first introduced in \cite{Saur_2013}, was inspired by the Jupiter-Io interaction and it works also in the case when the planet has no intrinsic magnetic field, but only a conducting atmosphere (ionosphere) or a conducting internal layer. In this scenario, magnetic power can propagate to the star only when the planet is inside the Alfvén surface of the stellar wind. The Alfvén surface is defined as the region of space where the Alfvén Mach number $M_{\rm A} \equiv v_{0}/v_{\rm A} < 1$, where $v_{0}$ is the velocity of the wind relative to the planetary body and $v_{\rm A} = B_{0} /\sqrt{\mu_{0} \rho}$ is the Alfvén velocity in the wind. In the latter term, $B_{0}$ is the intensity of the magnetic field in the wind, $\mu_{0}$ is the magnetic permeability of the vacuum, and $\rho$ is the wind plasma density. \\

The excited Alfvén waves move along the so-called Alfvén characteristic lines (cf. Sect.~2.1.1 of \cite{Saur_2013}) and, to observe an effect in the stellar atmosphere, at least one of these characteristic lines should connect the planet with the star. The power $P_{\rm SPI}^{\rm AW}$ that reaches the star can be computed numerically. When the Alfvén Mach number is small, it is possible to use the analytic approximation given in Eq.~(55) of \cite{Saur_2013}, that is
\begin{equation}
    P_{\rm SPI}^{\rm AW} = 2\pi R_{\rm eff}^{2}
    \frac{ \left( \overline{\alpha} M_{\rm A} B_{0} \cos \Theta \right)^{2}}{\mu_{0}} v_{\rm A},
\end{equation}
where $0 < \overline{\alpha} <1$ is a parameter that measures the interaction strength between the planet and the wind, $R_{\rm eff}$ is the effective radius of the planet that acts as an obstacle to the wind flow, $\Theta$ is the angle between the velocity of the stellar wind flow and the normal to the magnetic field of the wind at the position of the planet. All the quantities ($M_{\rm A}, B_{0}, \Theta, v_{\rm A}$) are to be evaluated at the position of the planet. The effective radius $R_{\rm eff}$ is equal to the radius of the planet $R_p$ when it has no intrinsic magnetic field, otherwise it depends on the magnetic field of the planet and can reach a maximum value of $\sqrt{3} R_{\rm obst}$, where $R_{\rm obst} = R_p (B_p/B_{0})^{1/3}$ with $B_p$ the magnetic field at the equator of the planet (cf. Sect.~2.3 of \cite{Saur_2013}). The equatorial field is half of the polar field in the case of a dipole field. \\

The maximum power available for SPI in the Alfvén wing model is obtained for $R_{\rm eff} = \sqrt{3} R_p (B_p/B_{0})^{1/3}$, $\overline{\alpha}=1$, $\Theta = 0$ and corresponds to
\begin{equation}
\begin{array}{lll}
P_{\rm SPI~(max)}^{\rm AW} & = & 6 \pi R_p^{2} (B_p/B_{0})^{2/3} M_{\rm A}^{2} B_{0}^{2} v_{\rm A}/\mu_{0} = \\
& = & (6\pi/\mu_{0}) M_{\rm A} R_p^{2} B_{0}^{4/3} B_p^{2/3} v_{0},
\end{array}
\label{p_aw_max}
\end{equation}
an expression that is rigorously valid in the limit $M_{\rm A} \rightarrow 0$.



\paragraph{The Magnetic Reconnection Model}

In this model, magnetic reconnection occurs between the stellar and planetary magnetic fields at the boundary of the planet's magnetosphere, requiring the planet to have an intrinsic magnetic field \cite{Lanza_2012}. The power released in this interaction is:

\begin{equation}\label{eq:Reconnection}
P^{\text{REC}}_{\text{SPI}} = \gamma \frac{\pi}{\mu_0} R^2 B_0^{4/3} B_p^{2/3} v_0,
\end{equation}
where $0 \leq \gamma \leq 1$ depends on the relative orientation of the magnetic field lines of the stellar and planetary fields, while the other symbols have already been introduced. The effective radius of the planetary magnetosphere is already included in Eq.~(\ref{eq:Reconnection}). The value of $\gamma \sim 1$ when the magnetic fields have opposite directions so that the reconnection has its maximum efficiency. By comparing Eqs.~(\ref{p_aw_max}) and~(\ref{eq:Reconnection}), we see that

\begin{equation}\label{p_ratio}
\frac{P^{\text{AW}}_{\text{SPI}}(\text{max})}{P^{\text{REC}}_{\text{SPI}}(\text{max})} = 6 M_A < 1,
\end{equation}
in the limit in which Eq.~(\ref{p_aw_max}) is valid. Equation~\ref{p_ratio} shows that the power made available by the Alfv\'en wing mechanism is usually smaller than that delivered by the reconnection mechanism provided that $M_{\rm A} \ll 1$. If one is interested in the order of magnitude of the available power, one can neglect the power conveyed by the Alfv\'en waves in comparison with that released by the reconnection. However, we notice that when the planet has no intrinsic magnetic field, $P_{\rm SPI}^{\rm REC} = 0$, only the power conveyed by the Alfv\'en waves can produce an SPI signal in the stellar atmosphere.


\paragraph{The Helicity Variation Model}

Ref.~\cite{Lanza_2012} considered the additional energy that can be released in a tall stellar magnetic loop that extends up to the orbit of the planet when we take into account the variation of the relative magnetic helicity\footnote{In the Sun, magnetic helicity variations have been invoked as triggers for large solar flares and coronal mass ejections (see \cite{Lanza_2012}, Sect.~2).} induced by the perturbations due to the orbital motion of the planet. The amount of available energy depends on the details of the stellar magnetic field configuration that is assumed to be force-free in the corona. The largest energy release is obtained when the field is a non-linear force-free field. \cite{Lanza_2012} considered a specific non-linear force-free field and found that the available power is at most $2-4$ times larger than that released by the magnetic reconnection as modeled above (see Sect. 3.3 in \cite{Lanza_2012}). Generalizing this result to a generic force-free field, we conclude that the helicity variation can provide an SPI power that is larger than that of the magnetic reconnection power, but of the same order of magnitude. Larger energy releases are occasionally possible when a stellar loop perturbed by the planet has stored a large amount of magnetic energy and is going to dissipate it producing a strong flare. In such a case, a small perturbation by the orbiting planet can trigger the energy release, thus producing a preference for flaring when the planet passes over an active region that is ready to flare (\cite{Lanza_2018}).

\paragraph{The Interconnecting Loop Model}

This model assumes that the magnetic field of the stellar corona can reconnect with that of the planet and form a loop whose field lines interconnect the surface of the star with that of the planet, assuming that its atmosphere is completely ionized (\cite{Lanza_2013}). The formation of such an interconnecting loop is possible only when the stellar field is potential and remains in such a state during the subsequent energy dissipation phase, which is possible provided that the energy dissipation is very efficient inside the loop. In this model, energy is continuously accumulated inside the loop because of the stress applied by the orbital motion at the planetary footpoint and is immediately released to keep the configuration of the field potential. Since a potential field is the minimum-energy configuration for assigned boundary conditions on the star and the planet, such a strong dissipation makes the field stable by maintaining it in the potential configuration.\\

The fractional planetary surface area \( f_{AP} \) that acts as the footpoint of the interconnecting loop is estimated from a model by \cite{Adams_2011}, which depends on the ratio \( \zeta = B_0/B_p \) between the stellar magnetic field \( B_0 \) at the planet’s location and the planetary magnetic field \( B_p \):

\begin{equation} \label{f_ap}
f_{AP} = 1 - \sqrt{1 - \frac{3 \zeta^{1/3}}{2 + \zeta}} \quad .
\end{equation}\\

The power in the interconnecting loop model is then:

\begin{equation} \label{p_il}
P^{\text{IL}}_{\text{SPI}} = \frac{2 \pi}{\mu_0} f_{AP} R^2 B_p^2 v_{\text{orb}},
\end{equation}
where \( v_{\text{orb}} \) is the planet’s orbital velocity. This expression assumes that the stellar rotation period is much longer than the planet's orbital period, so the loop footpoint velocity is effectively equal to \( v_{\text{orb}} \). Note that \( B_0 \) is indirectly included through \( \zeta \) in \( f_{AP} \).



\subsubsection{Application to the GJ 436 system}\label{subsub:Bp}
Not all the models described above can reproduce the observed excess power. In this section, we investigate which of the above SPI mechanisms can provide enough power to account for the observations. In particular, we will focus our calculations on the excess power observed in C-16, as we know that the planet was orbiting within the Alfvén surface at that time \cite{Vidotto_2023}. \\

The stellar and planetary parameters as well as the model of the stellar magnetic field are taken from \cite{Vidotto_2023}. We start by computing the power made available by the magnetic reconnection model and assume the following parameters in Eq.~\ref{eq:Reconnection}:
\begin{equation}
\begin{array}{ll}
 \gamma =1 \\
 R_p = 3. 85~R_{\oplus} = 24\,530~\mbox{km},\\
 B_{0} = 0.013~\mbox{G}, \\
 B_p = 1~\mbox{G}, \\
 v_{0} = 770~\mbox{km/s},
\end{array}
\end{equation}
where the planet radius comes from Table~1 of  \cite{Vidotto_2023}, and the magnetic field at the distance of the planet $B_{0}$ and the relative velocity between the planet and the wind $v_{0}$ come from their Model I (see Sect. 3 of  \cite{Vidotto_2023}). We adopt the maximum value of $v_{0}$ in their model to maximize the SPI power. For this model, we first need to assess whether it is capable of reproducing the observed powers; thus, the planetary magnetic field strength is assumed. The power made available according to Eq.~\ref{eq:Reconnection} is $P_{\rm SPI}^{\rm REC} = 3.5 \times 10^{16}$~W. This value should be regarded as an upper limit because we adopted the maximum value of $v_{0}$ from the considered wind model and $\gamma=1$. If we assume a stronger planetary field, the power made available by magnetic reconnection scales with $B_p^{2/3}$. For example, for a field of 10~G, we have a power of $1.6 \times 10^{17}$~W. To match the observations, we need $B_p \sim 3.1 \times 10^{4}$~G, that is, an unreasonably strong planetary field.\\

As described above, the power made available by the Alfvén wing mechanism is smaller than the one by the magnetic reconnection model ($P_{\rm SPI}^{\rm REC}$) by a factor of $6\,M_{\rm A}$. In our case, adopting Model I of \cite{Vidotto_2023}, the planet is close to the Alfvén surface where $M_{\rm A} = 1$, therefore, the validity of Eq.~\ref{p_aw_max} is questionable and the Alfvén wing power should be computed using a numerical method as detailed in \cite{Saur_2013}. However, that power is in the best case of the same order of magnitude as the reconnection power. Therefore, also the Alfvén wing model cannot account for the observation unless we assume an unreasonably strong planetary field. A similar conclusion is reached for the magnetic helicity variation model because even in that case the expected gain with respect to the reconnection model is of a factor $3-4$ in the best case assuming the same planetary magnetic field. \\

Lastly, we solve the equation of the interconnecting loop model for $B_p$. The power obtained in Sec. \ref{Sec: powers} for the C-16 observations is $P_{\text{C-16}} = 3.14 \times 10^{19}\,\text{W}$. The orbital velocity, assuming for simplicity a circular orbit, is $v_{\rm orb} = 2\pi a/P_{\rm orb} = 115.2$\,km/s. We consider two stellar magnetic fields at the distance of the planet of 0.013~G and 0.026~G, corresponding to Model I and II in \cite{Vidotto_2023}, respectively. Substituting into Eq.~\ref{p_il} and solving for $B_p$, we find a minimum planetary magnetic field from the interconnecting loop model of $B_p^{\rm V23-MI} = 10.36$~G and $B_p^{\rm V23-MII} = 9$~G. \\

The stellar magnetic field at the distance of the planet is a crucial parameter in the interconnecting loop model. Instead of using the model by  \cite{Vidotto_2023} (App.~B), we consider the most favorable case, that is, a radial stellar field with an intensity $B_{\rm star}=27$~G at the poles of the star as found using Zeeman Doppler Imaging in \cite{Bellotti_2023} following \cite{Lanza_2018}. Such a field scales with the distance $r$ as $B_{0}(r)= B_{\rm star} (r/R_{*})^{-2}$, where $R_{*}$ is the radius of the star. In our system, the planet mean distance is $a/R_{*}= 14.56$ (see Table~1 of  \cite{Vidotto_2023}), giving $B_{0}= 0.13$~G. Adopting such a value for the intensity of the stellar field at the distance of the planet, we find a minimum planetary magnetic field in the interconnecting loop of $B_p^{\rm L18} = 6.3$~G.  \\

These values represent lower limits for the planetary magnetic field because we are only detecting a fraction of the total energy dissipated in the SPI event, as discussed in Sec. \ref{Sec:Frac_Power}. Table \ref{Tab:Bfields} summarizes the resulting $B_p$ for various fractions of the detected power for each model of the stellar magnetic field in the interconnecting loop scenario.


\subsubsection{Planetary magnetospheric radius} \label{subsub:Rm}
With the information derived above, we can deduce the planetary magnetospheric radius $r_M$. This value refers to the front-facing radius of the magnetosphere --- the part pointing toward the star --- as it elongates in the opposite direction. Determining the value of $r_M$ is not straightforward, as it depends on various factors, such as the magnetic field generating it, $B_p$, and the surrounding environment. For the same $B_p$, if the planet is close to its star, its magnetosphere will be more compressed compared to a planet farther away due to a combination of the thermal pressure of the plasma, the ram pressure of the stellar wind, and the star's magnetic pressure (\cite{Ip_2004}, \cite{Vidotto_2013}). If the external pressure is sufficiently high, the planet’s magnetosphere could be compressed to the point that part of its atmosphere becomes exposed to stellar wind erosion, potentially leading to significant atmospheric loss (\cite{Vidotto_2013}). Some authors have provided analytical functions to estimate the approximate magnetospheric radius (\cite{Lanza_2009}, \cite{Vidotto_2013}, \cite{Lanza_2018}). Ref. \cite{Lanza_2018} derives the following expression for the magnetospheric radius:

\begin{equation}\label{Eq:rM_Lanza}
    r_M^{\rm L18} = 2(2f_0)^{1/3} \left[ \frac{B_p}{B(\textbf{r}_p)} \right]^{1/3} R_p \quad .
\end{equation}
where $f_0$ is a shape factor equal to 1.16 \cite{Lanza_2018}  and $\textbf{r}_p$ the position vector of the planet. Table~\ref{Tab:Bfields} reports the magnetospheric radius of GJ~436~b $r_M^{\rm L18} = 9 - 21\,R_p$ using this model assuming different fractions of the total power emitted in the calcium lines as discussed in Sec.~\ref{Sec:Frac_Power}. \\

However, these estimates are upper limits, as they primarily account for the stellar magnetic pressure, neglecting other pressures, such as ram pressure. A more detailed estimate of the planetary magnetosphere requires additional parameters of the planetary environment such as the stellar magnetic field and wind characteristics. Fortunately, due to the interest in GJ 436, previous studies have characterized the stellar environment to study wind-planet interactions (\cite{Bellotti_2023}, \cite{Vidotto_2023}). From the equilibrium between the planetary magnetic pressure and the total stellar pressure, \cite{Vidotto_2023} derived an analytical expression for the magnetospheric radius:

\begin{equation}\label{Eq:rM_Vidotto}
    r_M^{\rm V23} \simeq R_p \, f \left( \frac{(B_p/2)^2 /(8\pi)}{p_{\text{tot}}(a_{\text{orb}})} \right)^{1/6} \quad ,
\end{equation}
where $f \simeq 2^{2/6}$ is a correction factor, $B_p$ is the planetary magnetic field, and $p_{\text{tot}}(a_{\text{orb}})$ is the average total pressure of the stellar wind at the planetary orbital distance ($8.2 \times 10^{-6} \, \text{dyn/cm}^{2}$ and $59 \times 10^{-6} \, \text{dyn/cm}^{2}$ for Models I and II, respectively in \cite{Vidotto_2023}).\\

In their work, using two different stellar wind models and relying on a previous stellar magnetic field determination via spectropolarimetry (\cite{Bellotti_2023}), they derived magnetospheric radii for the planet GJ~436~b of $r_M = 5.2\,R_p$ and $r_M = 3.7\,R_p$ for each model. This value was estimated assuming a planetary magnetic field of $B_p = 2\, \text{G}$ and an SPI power emission via Alfvén wings of $P_{\text{SPI}}^{\rm AW} = 1.2 \times 10^{15} \, \text{W}$ and $P_{\text{SPI}}^{\rm AW} = 4.8 \times 10^{15} \, \text{W}$, respectively. In our case, since we have detected the magnetic interaction between the planet and the star, we can directly use our estimation of $B_p$. By substituting all the variables into Eq.~(\ref{Eq:rM_Vidotto}) we can calculate $r_M$ for the various values of $B_p$ in Table~\ref{Tab:Bfields}. The results are shown in Table~\ref{Tab:Bfields} ranging from $r_M^{\rm V23} = 6-20\,R_p$ for different percentages of the expected fraction of detected power, consistent with the \cite{Lanza_2018} model.




\begin{figure}[]
\centering
\includegraphics[width=\columnwidth]{ALL_AI_H08.png}
\caption{GLS periodogram for the same H-08 sub-epoch as shown in Fig.~\ref{Fig_1}, including all activity indicators available from HARPS. The last row represents the periodogram of the window function. The three horizontal red lines mark the 10, 1 and 0.1\% FAP levels from the GLS, while the gray vertical ones mark the signals at 2.81\,d, 2.64\,d and 2.46\,d. }
\label{HARPS_2008_All}
\end{figure}

\begin{figure}
    \centering
    \includegraphics[width=\columnwidth]{ALL_AI_H08-Entire-epoch.pdf}
    \caption{Same as Fig.~\ref{HARPS_2008_All}, but covering the entire 2008 observational season.}
    \label{HARPS_2008full_All}
\end{figure}

\begin{figure}
    \centering
    \includegraphics[width=\columnwidth]{ALL_AI_H09.pdf}
    \caption{Same as Fig.~\ref{HARPS_2008full_All}, but covering the entire 2009 observational season.}
    \label{HARPS_2009_All}
\end{figure}

\begin{figure}
    \centering
    \includegraphics[width=\columnwidth]{ALL_AI_H10.pdf}
    \caption{Same as Fig.~\ref{HARPS_2008full_All}, but covering the entire 2010 observational season.}
    \label{HARPS_2010_All}
\end{figure}

\begin{figure}
    \centering
    \includegraphics[width=\columnwidth]{ALL_AI_H20.pdf}
    \caption{Same as Fig.~\ref{HARPS_2008full_All}, but covering the entire 2020 observational season.}
    \label{HARPS_2020_All}
\end{figure}

\begin{figure}[]
\centering
\includegraphics[width=\columnwidth]{ALL_AI_C16.png}
\caption{GLS periodogram for the same C-16 sub-epoch as shown in Fig.~\ref{Fig_1}, including all activity indicators available from CARMENES. The last row represents the periodogram of the window function. The three horizontal red lines mark the 10, 1 and 0.1\% FAP levels from the GLS, while the gray vertical ones mark the signals at 2.81\,d, 2.64\,d and 2.46\,d. }
\label{CARMENES_2016_All}
\end{figure}

\begin{figure}[]
\centering
\includegraphics[width=\columnwidth]{ALL_AI_C24.png}
\caption{Same as Fig.~\ref{CARMENES_2016_All}, but for C-24.}
\label{CARMENES_2024_All}
\end{figure}


\begin{figure}
\centering
\includegraphics[width=\columnwidth]{Final_RP_HARPS_ALL.png}
\caption{\textbf{Rolling periodogram of the $I_{\rm CaII}$ from all the combined observations taken with HARPS  between 2006 and 2020.} The rows, from the top to the bottom, are the periodogram of each consecutive sub-epoch containing m\,=\,40 datapoints. On the top the GLS periodogram of all the combined observations as a reference. Three vertical white lines mark the signals at 2.81\,d, 2.64\,d and 2.46\,d. Three other black lines mark the same periods in the top GLS periodogram.}
\label{fig:Rolling_periodogram_HARPS}
\end{figure}

\begin{figure}
\centering
\includegraphics[width=\columnwidth]{Final_RP_CARMENES16.png}
\caption{\textbf{Same as Fig.~\ref{fig:Rolling_periodogram_HARPS}, but for the pEW'(CaII~{\textsc{IRT-a}}) covering the entire 2016 epoch observed with CARMENES.}}
\label{fig:Rolling_periodogram_CARMENES}
\end{figure}




\bibliography{sn-bibliography}


\end{document}
===== END FILE: article.tex =====
